\chapter{Introduction}
\label{chap:introduction}

In recent years, space communications and satellite networks have undergone a monumental paradigm shift. The rapid miniaturization of space components, combined with drastically reduced launch costs, has catalyzed the deployment of Mega-Constellations operating in Low Earth Orbit (LEO)~\cite{survey_orbital_edge_2023}. Traditionally, satellite architectures functioned strictly under a "bent-pipe" model, where satellites served as passive transparent relays forwarding raw telemetry or communication signals back to Ground Stations (GS) for processing. However, the unprecedented volume of data produced by modern spaceborne applications—such as high-resolution Earth Observation (EO)~\cite{leyva2023satellite}, environmental monitoring, and ubiquitous Internet of Things (IoT) sensing~\cite{tiansuan2023}—has rendered this traditional pipeline highly inefficient. 

To overcome the severe bandwidth limits of space-to-ground downlinks and the high propagation delays associated with remote terrestrial processing, the paradigm of \textbf{Orbital Edge Computing (OEC)}—also broadly termed Satellite Computing (SC)—has emerged~\cite{denby2019orbital, denby2020asplos}. By equipping LEO satellites with onboard processing hardware, data can be computed directly in space close to the source, transmitting only actionable insights or highly compressed results to ground users. Despite its immense potential, executing compute-intensive tasks on resource-constrained satellite nodes introduces complex orchestration challenges that require advanced decision-making frameworks~\cite{survey_orbital_edge_2023, icdcs_satellite_2025}.

\section{Motivation and Context}
\label{sec:intro_motivation}

The transition towards Orbital Edge Computing is driven by two main technological trends: the exponential increase in space-generated data and the strict temporal requirements of next-generation applications. Applications such as real-time disaster response, autonomous maritime navigation, and defense surveillance demand sub-second latency for task execution and reporting~\cite{leyva2023satellite}. Processing these tasks on Earth introduces unavoidable latency bottlenecks due to intermittent Ground Station contact windows, dynamic weather interference, and severe downlink congestion.

LEO megaconstellations (such as Starlink, OneWeb, and Kuiper) consist of hundreds to thousands of satellites orbiting at altitudes between 300 km and 2000 km. These satellites move at high orbital velocities (approximately $7.5 \text{ km/s}$ relative to the Earth's surface), resulting in a continuously changing network topology. Inter-Satellite Links (ISLs)—utilizing high-throughput Free-Space Optics (FSO) or Radio Frequency (RF) channels—enable direct communication between adjacent satellites without routing signals down to the ground~\cite{chen2021isl, chaudhry2021laser}. 

Consequently, a LEO constellation effectively forms a distributed, high-velocity edge computing network in orbit. However, executing applications across this orbital fabric requires intelligent \textbf{Task Offloading and Routing Orchestration}~\cite{zhong2025leo, fan2025satellite}. Ground users or source satellites generate computational tasks that must be dynamically assigned to suitable processing nodes across the constellation. Deciding \textit{where} and \textit{when} to execute a given task, while routing its associated payload across multiple dynamic ISL hops, constitutes a high-dimensional optimization problem that directly governs overall network performance.

\section{Problem Statement}
\label{sec:intro_problem_statement}

Orchestrating task offloading in Orbital Edge Computing environments presents a unique set of constraints and operational challenges that set it apart from conventional terrestrial cloud/edge systems~\cite{icdcs_satellite_2025}:

\begin{itemize}
    \item \textbf{Severe Energy Budget Limits:} Satellite Edge Nodes (SENs) rely on solar panels and onboard battery storage. Computationally intensive task execution and RF/optical transmissions consume substantial power. Over-utilizing specific nodes leads to rapid battery depletion, thermal overhead, or total operational failure~\cite{icdcs_satellite_2025}.
    \item \textbf{Heterogeneous and Limited Hardware Resources:} Onboard processing units (e.g., embedded CPUs, specialized AI accelerators) are constrained by strict Space, Weight, and Power (SWaP) requirements, as well as radiation-hardening constraints~\cite{denby2020asplos}.
    \item \textbf{Dynamic Topology and Intermittent Connectivity:} High orbital mobility leads to frequent topology reconfigurations, changing link distances, and periodic loss of line-of-sight. Task routing paths must be continuously recalculated to avoid dropped links~\cite{chen2021isl}.
    \item \textbf{Strict Temporal Constraints and Queuing Delays:} Tasks possess strict deadlines ($D_r$). If the sum of transmission latency, queuing delay, and computation time exceeds $D_r$, the task fails, resulting in wasted system energy and degraded quality of service.
\end{itemize}

To achieve theoretical optimality in task placement and energy management, mathematical optimization tools such as \textbf{Integer Linear Programming (ILP)} are commonly employed~\cite{wolsey1998integer, boyd2004convex}. An ILP formulation can search the joint solution space of task routing, CPU allocation, and energy minimization to find globally optimal schedules. 

However, centralized ILP solvers suffer from inherent computational complexity (NP-hard)~\cite{wolsey1998integer}. In dynamic satellite networks with continuously arriving workloads, solving an ILP instance for every individual task request introduces unacceptable computational overhead and processing latency at the central orchestrator. Conversely, accumulating task requests to solve them in batches introduces a trade-off: larger batches allow the ILP solver to achieve superior global resource allocation, but prolonged waiting times in orchestrator queues increase the risk of task deadline expiration (\textit{Deadline Exceeded} failures). 

\section{Objectives and Main Contributions}
\label{sec:intro_objectives}

The main objective of this thesis is to design, implement, and rigorously evaluate a \textbf{Centralized ILP Orchestration Framework} optimized for task routing and resource allocation in LEO satellite networks, incorporating a \textbf{Dynamic Batching Mechanism} to balance solver computational complexity with network queuing delay.

Specifically, this work addresses the key question: \textit{To what extent can a centralized optimization approach maintain global allocation efficiency in high-density satellite networks before queuing and solver overhead cause systemic performance degradation?}

The main contributions of this thesis are summarized as follows:

\begin{enumerate}
    \item \textbf{Formalization of the Centralized ILP Optimization Model:} We formulate a comprehensive ILP mathematical model that jointly optimizes task routing across Inter-Satellite Links and computational resource allocation on Satellite Edge Nodes. The model supports dynamic weighting between latency optimization (\textit{Mapping Time}) and total energy preservation (\textit{Mapping Energy}).
    \item \textbf{Design of an Integrated Dynamic Batching \& Timeout Mechanism:} To mitigate the computational overhead of solving the ILP model continuously, we integrate an adaptive batching policy governed by Batch Size ($BS$) and Batch Timeout ($BT$). Furthermore, we propose a mathematical deadline compensation formula to adjust task deadlines based on orchestrator queuing duration ($T_{batch}$).
    \item \textbf{Design and Integration of the Centralized Orchestration Layer:} We design and implement the centralized ILP orchestration module and dynamic batching logic, integrating them into an existing Python/\texttt{SimPy} discrete-event simulation framework~\cite{simpy2023, pulp2011}. The resulting environment accurately captures request arrival dynamics ($\lambda$), inter-satellite communication links, queue management, hardware processing, and fleet energy consumption under heterogeneous workload conditions.
    \item \textbf{In-Depth Performance and Bottleneck Analysis:} We perform extensive stress testing under high workload arrival rates to identify system collapse thresholds, distinguishing between physical resource saturation (CPU/Network capacity) and temporal violations (\textit{Deadline Exceeded}). We also conduct sensitivity analyses regarding deadline relaxation ($\Delta D$).
    \item \textbf{Comparative Benchmark Against State-of-the-Art Solutions:} We evaluate the centralized ILP orchestrator against established baseline heuristics (\textit{OrbitAware}, \textit{DTS-base}, \textit{DTS-APopt}) and distributed ILP formulations from the literature~\cite{icdcs_satellite_2025}, establishing clear operational trade-offs regarding task completion rate, overall latency, and energy efficiency.
\end{enumerate}

\section{Thesis Outline}
\label{sec:intro_outline}

The remainder of this thesis is structured as follows:

\begin{itemize}
    \item \textbf{Chapter 2 (Background and State of the Art):} Reviews the fundamental concepts of Satellite Computing, LEO network dynamics, and existing literature on task offloading. It provides a taxonomy of centralized versus distributed orchestration schemes and details key reference algorithms.
    \item \textbf{Chapter 3 (System Architecture and ILP Formulation):} Presents the formal mathematical system model, decision variables, objective functions (\textit{Mapping Time} vs. \textit{Mapping Energy}), constraints, and the dynamic batching/timeout mechanisms.
    \item \textbf{Chapter 4 (Performance Analysis and Evaluation):} Describes the \texttt{SimPy} discrete-event experimental setup and evaluates the proposed centralized orchestrator under varying batching configurations, heavy-load conditions, and deadline relaxation scenarios.
    \item \textbf{Chapter 5 (Comparative Benchmark):} Details the comparative analysis benchmarking our centralized orchestrator against state-of-the-art distributed heuristics and formulations, highlighting key trade-offs in latencies, energy consumption, and global optimality.
    \item \textbf{Chapter 6 (Conclusions and Future Work):} Summarizes the main findings of this research, outlines limitations of centralized orbital orchestration, and proposes future research directions toward hybrid adaptive architectures.
\end{itemize}